\section{Compatibility between the two units}\label{sec:phases}

A witness must respect the images already fixed in its leaves. This
requires four bits, each obtained by evaluating a functional from one
unit on a vector from the other. We estimate the probability that the
bits have the prescribed values.

A marked unit carries profiles \(\lambda_i\in C_i\), of bounded total
rank. Put
\begin{equation}\label{eq:delta-phase}
 \Delta_A=U_{A_1}\lambda_1+U_{A_2}\lambda_2,
 \qquad u_{A,S}=\sum_{i\in S}u_{A_i}.
\end{equation}
For independent units \(A,B\), the relevant bits are
\[
 u_{B_1}(\Delta_A),\quad u_{B_2}(\Delta_A),\quad
 u_{A_1}(\Delta_B),\quad u_{A_2}(\Delta_B).
\]
For example, if \(p_i=T(\lambda_i,\cdot)+a\), the balance equation
\eqref{eq:recipe-balance} becomes
\(u_{B_j}(\Delta_A)=a(\lambda_1)+a(\lambda_2)\).
The other prediction classes lead to parities of the same four bits.

\subsection{A cover alternative}

A cover is a collection of ambient subspaces \(C_e\). It covers a symmetric
profile \(x\) when
\[
 x_e\in C_e\otimes V+V\otimes C_e\quad\text{for every }e.
\]
Its dimension is \(\sum_e\dim C_e\).

% lax begin Lax793645.PhaseAlternative
\begin{proposition}\label{thm:phase-alternative}
Suppose \(\Delta_A\ne0\) and its total rank is bounded. From a law with
the prepared marginal and joint density bounds, one can retain a fixed
positive fraction on which one of the following holds.

Either every nonconstant character of the four displayed bits has mean
of absolute value less than \(.005\), or all \(\Delta_A\) are covered by
one cover of bounded dimension and every character using a single
endpoint on at least one side has mean at most \(2^{-.24N}\).
The bounds needed for the cover and for the retained mass precede the
choice of the channel dimension.
\end{proposition}
% lax end

We give the estimates behind the proposition. They also specify which
uses of independence are legitimate.

% lax begin Lax793645.RankOfIndependentSums
\begin{lemma}\label{lem:no-cover-rank}
Let independent tensors \(T_1,\ldots,T_s\) have total rank at most \(r\).
There is \(q_r>0\), depending only on \(r\), such that if every cover
of dimension at most \(2rs\) has probability less than \(p\), then
\[
 \PP\left(\rank\sum_jT_j<q_rs/4\right)\le4^sp^{q_rs/2}.
\]
It suffices to take \(q_r=2^{-100r-10}\), with \(s\) a sufficiently
large power of two.
\end{lemma}
% lax end
% lax begin Lax793645Proofs.Manuscript.RankOfIndependentSums.no_cover_rank
\begin{proof}
Expose some indices and quotient by their column and row spans. Among
\(b\) remaining indices, let \(\delta_C,\delta_R\) be the deficits
between the sum of their individual projected dimensions and the dimension
of their combined span. The rank of the projected sum is at least
\[
 \#\{j:\overline T_j\ne0\}-\delta_C-\delta_R.
\]
Place the projected tensors on separate direct summands first. The maps
adding the summands have kernels of dimensions \(\delta_C,\delta_R\);
applying them loses at most those ranks.

It remains to choose which indices to expose. Write \(D(E)\) for the
sum of the dimensions of the row and column spans of the tensors indexed
by \(E\). Put \(H_b=\sum_{j=1}^b1/j\). Among nonempty sets of unexposed
indices \(R\), minimize
\[
  u(R)-H_b/4,\qquad b=|R|,\quad E=[s]\setminus R,\quad
  u(R)=\frac{D([s])-D(E)}{b}.
\]
Initially \(u([s])\le 2r\), and \(u(R)\ge0\) for every \(R\).
Each doubling of \(b\) adds at least \(1/2\) to \(H_b\). Thus
\(b\le q_rs\) would give \(H_s-H_b\ge(100r+10)/2>8r\),
contrary to minimality. In particular \(b>1\).

For \(i\in R\), let \(t_i=D(E\cup\{i\})-D(E)\). Removing \(i\) from
\(R\) changes its average from \(u(R)\) to \((bu(R)-t_i)/(b-1)\).
Minimality gives \(t_i\le u(R)+(b-1)/(4b)\). Summing over \(i\)
gives \(\sum_{i\in R}t_i-(D([s])-D(E))\le b/4\).
This bounds the combined row and column deficit after projection.

If the original sum has rank less than \(q_rs/4\), at least \(q_rs/2\)
remaining tensors vanish in this quotient. They lie in the cover generated
by the exposed tensors. For fixed disjoint exposed and vanishing index
sets, condition only on the exposed tensors. The rest are still independent,
so this event has probability at most \(p^{q_rs/2}\). There are at most
\(4^s\) choices of the two index sets. For symmetric tensors, combine the
row and column cover spaces into the single cover stated above.
\end{proof}
% lax end

% lax begin Lax793645Proofs.Manuscript.PhaseAlternative.one_phase_alternative
\begin{proof}[Proof of Proposition~\ref{thm:phase-alternative}]
First take \(s\) proportional to \(N\), with a sufficiently small fixed
coefficient. If no cover of dimension \(2rs\) has appreciable mass,
expand an even \(s\)-th moment of a phase mean. For independent free
channel columns, a single channel pair has expectation
\(2^{-\rank W}\) against the matrix \(W\). Hence the moment is at most
\[
 \E 2^{-h\rank(\Delta_{A_1}+\cdots+\Delta_{A_s})}
 \le4^sp^{q_rs/2}+2^{-hq_rs/4}.
\]
The actual raw channel law is bounded by a fixed multiple of this free
law. This follows from the quadratic frame counts, after summing over the
unspecified Gram data. That multiple depends on the fixed channel size,
but not on \(n\).
Choose \(p\) small and then \(h\) large. Markov's inequality beats the
joint density cap and the \(2^{O(rN)}\) possible marked matrices. The
extra phase from the opposite mark can be fixed before this calculation
and then included in that finite union. Every nonconstant four-bit
character is therefore small.

Otherwise restrict to a cover of dimension \(O(s)\) carrying positive
mass. The individual primal images avoid this fixed cover with
exponentially small exception. Indeed a uniform image of dimension
\(\dim B\) meets a fixed space of dimension \(d\) with probability
at most \(2^{\dim B+d-N+O(1)}\). The marginal cap preserves this estimate.
Pass to the quotient by the cover and lift a basis of the common image
of the two marked tensors. Their difference has the symmetric form
\begin{equation}\label{eq:symmetric-offsets}
 \Delta_e=\Omega_eZ_e\trans+Z_e\Omega_e\trans+
                \Omega_eH_e\Omega_e\trans.
\end{equation}
The columns of \(\Omega_e\) belong to the cover and record the differences
of the two lifts. The companion columns \(Z_e\) can be taken in either
endpoint's primal image. After removing redundant offsets, their total
number \(t\) lies between 1 and a constant depending on \(r\).
The offset vectors and coefficient matrices form a record with at most
\(2^{.002N}\) possibilities, by making the coefficient in \(s=O(N)\)
small enough.

For an ambient quadratic form \(q_u\), write
\(\Phi_u(v)=\beta_{q_u}(v,\cdot)\), identifying the dual with the
ambient coordinate space. Evaluating \eqref{eq:symmetric-offsets} gives
\(Z\cdot\Phi_u(\Omega)\) plus terms determined by the offset record.
The values \(q_u(\omega)\) on its basis must also be recorded; polar
values alone do not determine the diagonal terms.

The companion tuple has point masses at most \(2^{-tN+o(N)}\).
For a single endpoint, if its \(X\)-evaluations on a fixed offset basis
are independent, its companion tuple together with \(\Phi_u\) on that
basis has point masses at most \(2^{-(t+t')N+o(N)}\).
To verify the latter count, condition on the primal and \(X\) columns.
The formula
\[
 \Phi_u(v)=\sum_j\bigl(\beta_0(X_j,v)\,\beta_0(Y_j,\cdot)
                  +\beta_0(Y_j,v)\,\beta_0(X_j,\cdot)\bigr)
\]
is linear in the remaining \(Y\) columns. Modulo the span of the \(X\)'s,
the independent evaluation coefficients give rank at least \(t'(N-h)\).
Include the companion image prescription and the bounded coefficient
record. Frame counting gives the asserted point bound. Failure of the
independence condition on a fixed \(t'\)-space has probability
\(O(2^{t'-h})\).

Apply \eqref{eq:walsh-bound} to the two companion/action tuples in a phase.
A singleton endpoint supplies the stronger point bound on one side, so
its contribution is exponentially small apart from the preceding rank
failure. Sum over the offset records. Their \(.004N\) cost leaves room
for the exponent \(.24N\).

There is one further choice. If a bounded-dimensional space contains the
offsets on a fixed positive mass, retain that mass and require both
endpoints to pass its channel rank test. This gives the second alternative
with no remaining fixed error in the singleton estimates.

If there is no such space, the full endpoint sum also has cancellation.
For completeness, fix an offset list and call an action value frequent
when its joint mass with a companion record exceeds
\(2^{-(t+.5)N}\). The companion bound puts at most \(2^{.51N}\) values
on its frequent list. If frequent pairs had appreciable mass, choose a
fixed number \(L\) of independent probes. At each step the absence of
a small cover supplies a fresh offset direction with fixed positive
probability. Fix a successful probe list by averaging. After a small
rank-test deletion, prescribing all its action values costs at least
\(.99LN\) bits, whereas the frequent lists allow only \(.51LN\) bits.
The raw event thus has probability at most
\[
 2^{(D+.01)N+O(n)-.48LN}.
\]
Choose \(L\) before \(h\), large enough to beat the fixed joint cap.
This contradicts the positive probe mass. Off the frequent event, the
Walsh estimate makes the full-sum phase small as well. Interchanging the
two units covers all characters. Every dimension, mass threshold, and
probe count used here is chosen before \(h\).
\end{proof}
% lax end

\subsection{The probability of a zero phase}

If the first phase alternative holds, Fourier inversion gives probability
at least \((1-15\cdot.005)/16>.05\) to every four-bit pattern.
For a set \(S\) of endpoints put
\[
 F_S(A,B)=\sum_{i\in S}u_{B_i}(\Delta_A).
\]
This matrix factors through the ambient space of \(g\) component
matrices, so its rank is at most \(gN^2\) for every unit law.

Let \(M\) have rank at most \(k\ge1\) over any field. Draw \(A,B\)
independently from the same law \(\mu\), and put \(p=\PP\{M(A,B)\ne0\}\).
If every diagonal entry is nonzero, then
\begin{equation}\label{eq:rank-compatibility}
 \PP\{M(A,B)\ne0\}\ge\frac{1}{4k}.
\end{equation}
More generally, for every \(t>0\) there is a set \(I\) of mass at least
\(1-p/t-kt\) with \(M=0\) on \(I\times I\). Discard the rows whose
support has mass at least \(t\). Their indices have mass at most \(p/t\).
Choose a basis from the remaining rows and discard the indices in their
supports, of total mass at most \(kt\). Every remaining row vanishes at
every remaining column. If the diagonal is nonzero, \(I\) must be empty.
Taking \(t=1/(2k)\) proves \eqref{eq:rank-compatibility}.

Write \(J\) for the all-one matrix. Over \(\F\), the support of
\(J+F_S+F_S\trans\) records the zeros of \(F_S+F_S\trans\). Its diagonal
is one and its rank is at most \(1+2gN^2\). Thus the reciprocal zero
parity has probability at least \(1/[4(1+2gN^2)]\). This covers equal
active predictions and a single active endpoint with \(\alpha_i=1\).

For the four prescribed bits take \(F=F_{\{1,2\}}\) and put \(r=gN^2\).
If \(F(A,A)=0\) for every \(A\), put \(Q=J+F\). The entrywise product
\(Q\circ Q\trans\) has diagonal one. Its support records precisely the
pairs for which both directed entries of \(F\) vanish. Multiplying
rank-one expansions gives \(\rank(Q\circ Q\trans)\le(r+1)^2\).
Hence the probability of two zeros is at least \(1/[4(r+1)^2]\).

The singleton characters are exponentially small. Fourier inversion on
the four bits therefore gives
\[
 \PP\{\text{all four bits equal }c\}
 =\tfrac14\PP\{F(A,B)=F(B,A)=0\}+O(2^{-.24N})
\]
for either prescribed \(c\). These bounds cover the compatibility
conditions listed in the next section. In particular the prescribed
four-bit pattern has mass at least \(1/[16(gN^2+1)^2]-2^{-.24N}\).
The ambient rank estimate applies to every unit law. The character error
uses the singleton phase estimate. For a single active endpoint, only
the character of its reciprocal parity is needed: if its bias has
absolute value at most \(0.1\), either prescribed bit has probability
at least \(0.4\).

There is a direct count for the all-zero target when both endpoints are
inactive. At either endpoint \(p_i(\lambda_i)=T(\lambda_i,\lambda_i)
+a(\lambda_i)=0\). Since also \(u_iU_i=0\), each listener vanishes on
its own unit's offset. Suppose the offsets span a space of dimension
\(d\). Each matrix \(L_i(A,B)=u_{B_i}(\Delta_A)\) then has rank at
most \(d\). Put \(Q=(J+L_1)\circ(J+L_2)\). The matrix
\(Q\circ Q\trans\) has ones on its diagonal and rank at most
\((d+1)^4\). Its support is exactly the event that all four readings
are zero. The row support bound gives this event probability at least
\(1/[4(d+1)^4]\). For \(d=O(N)\) this is enough for the rare case.

For a prescribed reciprocal one, use \(L_i+L_i\trans\). For four
prescribed ones, use \(L_1\circ L_2\circ L_1\trans\circ L_2\trans\).
Each matrix has rank at most \(k=(d+1)^4\). Taking \(t=1/(4k)\) above
shows that either compatibility has mass at least \(1/(16k)\), or a
set of mass greater than \(1/2\) contains no compatible pair. Thus, on
a span within the label budget below, it suffices to find a compatible
pair in every such set. The phase alternative handles the remaining cases.

Suppose all retained offsets lie in a space \(V\) of dimension \(d\).
Label a unit \(A\) by \(\Delta_A\in V\) and the restrictions of
\(u_{A_1},u_{A_2}\) to \(V\). Each restriction is a linear functional
on the space of symmetric matrices. The label takes at most \(2^{3d}\)
values. Two labels determine all four readings, since each reading
evaluates one stored functional on the opposite offset. Split each leaf
by its label and repeat peeling. All descendants retain their label.
Compatibility is now constant on each pair of leaves, and there are at
most \(2^{3d}\) possible sets of accepting opposite units. A small
geometric cover supplies such a space with \(d\le C_sN\).

\subsection{Safety on compatible pairs}

% lax begin Lax793645.SafeCompatiblePairs
\begin{lemma}\label{lem:injection}
For a pair event \(E\) whose accepting sets have the preceding bound,
the preparation can ensure
\[
 \PP(E\text{ and an unsafe table})
 \le .0005\,\PP(E)+2^{-.005N}.
\]
In particular this loss is small relative to inverse-polynomial
compatibility.
\end{lemma}
% lax end
% lax begin Lax793645Proofs.Manuscript.SafeCompatiblePairs.coded_safety
\begin{proof}
Consider a fixed endpoint and component. Failure means that a nonzero
pinned primal image has channel evaluation in the opposite stored channel
space. Fix an accepting set of mass at least \(2^{-.01N}\).
Under its conditional law, draw \(\ell\) independent primal pin spaces,
where \(\ell\) is a small fixed multiple of \(N/(K+1)\).
The whole primal images are independent of the span of earlier witnesses
except with exponentially small probability, by the marginal frame bound.
If the conditional failure rate exceeds \(\eta\), all \(\ell\) tests
therefore succeed with probability at least \((\eta/2)^\ell\).

Conversely, fix independent witnesses. There are at most \(2^{K\ell}\)
ways to select them from the pins. Their channel evaluations are uniform,
up to the frame-count factor. Requiring them all to lie in a stored space
of dimension at most \(K\) costs \(2^{-(h-K)\ell+O(1)}\).
The possible stored channel spaces contribute only a fixed factor.
Choose \(h\) so this upper bound, after paying the marginal density cap
and the lower test probability, beats \(2^{C_sN}\).
A union over the accepting family then leaves an exponentially small
exception. Accepting sets of mass below \(2^{-.01N}\) already have that
small total pair contribution. Sum over the bounded number of endpoint
and component tests, choosing \(\eta\) accordingly. The conservative
parameter margins give the displayed constants.
\end{proof}
% lax end
